% ---------------------------------------------------------------------
%  04-model.tex -- the HPB model in plain language
%  The precise definitions, rules R1--R7 and properties live in
%  appendix-a-formal.tex (labels def:*, rule:*, prop:*).
% ---------------------------------------------------------------------
\section{The HPB Model}
\label{sec:model}

HPB has very few concepts, on purpose. Every concept had to earn its place
by answering a question that I was asked in real product reviews and could
not answer from the tools I had: \emph{what did this version contain,
where was this defect seen, what is still open, and where did this
reference come from?} This section explains the concepts in plain
language, with the symbols introduced only where later sections use them.
The precise definitions, the rules R1--R7 in their exact form and the
properties that follow from them are collected in \cref{app:formal}; a
tool builder should read that appendix, a product manager can stop here.

\subsection{Product lines and baselines}
\label{sec:model-lines}

\paragraph{Product line.}
A \emph{product line} $P$ is one product reference followed over time: a
device model with its manufacturer reference, a firmware, an application
or a service (\cref{def:line}). This is narrower than the ``product line''
of software product line engineering, which means a family of variants; in
HPB a line is \emph{one} reference and its history. The portfolio is the
set of all lines, and lines are related by descent: when a hardware
revision creates a new reference, the new line records the old one as its
parent. The result is a family tree of references, the \emph{line
forest}.

\paragraph{Baseline.}
A \emph{baseline} $v$ is a named, identifiable state of a line:
\texttt{v2.3.1}, \texttt{Rev~B}, \texttt{PC-7} (\cref{def:baseline}). I
use the word in its configuration-management sense, a reference
configuration against which change is
defined~\citep{ieee828,iso10007}. Every baseline has
\begin{itemize}[leftmargin=1.5em,itemsep=1pt,topsep=2pt]
  \item a \emph{kind}: firmware, software, hardware, or \emph{other}, the
        last one used for \emph{process baselines}, i.e., a change to a
        manufacturing process, test procedure or supplier that changes the
        product as built without a new design revision;
  \item a \emph{status}: \emph{planned} (a target with a date) or
        \emph{released} (a fact with a date);
  \item a name that is unique within its line. Any versioning convention
        works, for example Semantic Versioning~\citep{semver}.
\end{itemize}
A status only moves from planned to released, and a released baseline is
immutable. That one constraint is what makes everything else in HPB
trustworthy: if a released version could still absorb work, no change log
and no defect count would ever be final.

\paragraph{What ``released'' means.}
Every time-based metric depends on it, so HPB fixes a default per kind.
Firmware and software are released when the build is published to
customers; hardware when its engineering change order is approved
\emph{and} its production effectivity is recorded, with the change-order
identifier; a process change when it becomes effective in production. You
may define it differently, but write your definition down.

\paragraph{New revision or new reference?}
This question comes back with every hardware revision, and HPB settles it
with the test that engineering change management already
uses~\citep{jarratt2011ec,stark2015plm}: \emph{form, fit and function}. A
revised unit that is not interchangeable with the previous one in form,
fit or function gets a new reference and starts a new line (a \emph{line
fork}, below). An interchangeable change, such as a component substitution
that keeps the reference, is simply a hardware baseline on the existing
line.

\paragraph{Release streams.}
Products are rarely maintained on a single sequence of versions. Each line
has a \emph{main stream} and, when needed, \emph{maintenance streams} that
leave a released baseline of another stream, for example a long-term
support branch of a firmware (\cref{def:mainline}). The ordered baselines
of the main stream form the \emph{main line} $M(P)$, the thick horizontal
stroke of every HPB diagram. Release dates order baselines only
\emph{within} a stream: a patch \texttt{v1.4.1} on a stream forked at
\texttt{v1.4} follows \texttt{v1.4}, even if it ships after \texttt{v2.0}
on the main stream, and it does not contain anything that only went into
\texttt{v2.0}. A stream never changes the product reference; that is what
line forks are for (\cref{fig:backport}).

\paragraph{Variants and configurations.}
Regional or certification variants, OEM builds and members of a product
family are \emph{sibling} lines: distinct references, each with its own
history. Work that applies to all of them belongs on the line of the shared
element, typically a firmware or software line, instead of being
duplicated. Each firmware or software baseline records the hardware
references it is released to run on (its \emph{applies-to} set), and a
\emph{configuration} is a combination of baselines that is actually
possible in the field, such as (\texttt{Rev~A}, firmware \texttt{v2.0},
app \texttt{5.1}) (\cref{def:config}). The applies-to relation is the
compatibility matrix of the portfolio.

For a single device with its own firmware, a \emph{combined} line with
both hardware revisions and firmware releases is simplest, and most
examples here use it. Switch to separate lines as soon as one firmware
build runs on several hardware references (\cref{fig:shared}).

\subsection{Branches}
\label{sec:model-branches}

A \emph{branch} $b$ is one unit of change intent on a product: an
improvement to make, a defect to fix, an engineering task to do, a lab
evolution to qualify (\cref{def:branch}). It is \emph{not} a code branch.
A single HPB branch may be implemented by many commits in several
repositories, by a laboratory campaign or by a hardware redesign. HPB
records where the intent started, where it ended, and what happened in
between, and nothing more.

\paragraph{Four types.}
Every branch has exactly one type $\tau(b)$, as in \cref{sec:onepage}:
$\mathsf{I}$ improvement (from customers, sales, competitor analysis or the
roadmap), $\mathsf{L}$ laboratory evolution (evaluation, qualification or
prototype work), $\mathsf{X}$ issue (a defect observed in a released
baseline) and $\mathsf{T}$ task (work that keeps the product operable, such
as a certification renewal or a last-time buy). In a portfolio with
constant rotation, the balance between making the product better
($\mathsf{I}$, $\mathsf{L}$) and keeping it correct ($\mathsf{X}$,
$\mathsf{T}$) is one of the first things management asks about, so the
notation puts the two groups on opposite sides of the main line.
Sub-tasks are not branches: they are \emph{steps} of the branch they
belong to, so that each piece of intent is counted once (\cref{rule:typed}).

\paragraph{The fork baseline.}
Every branch starts from exactly one released baseline, its \emph{fork
baseline} $f(b)$, fixed when the request is accepted at triage
(\cref{rule:origin}). For an issue, it is the version in which the defect
was \emph{observed}, as reported. For the other types, it is the version
the change is specified against, by default the latest release of the main
stream. When a defect was reported against a combination of baselines,
the full observed configuration is kept and one of its baselines is the
primary one, on whose line the branch is drawn.

The fork baseline records what was known at intake and is never
rewritten by what is learned later: other affected versions go to an
\emph{also affects} list, and the version that \emph{introduced} the
defect is recorded separately. I insist on this because the first thing
lost when tickets are edited over months is the one fact quality
reporting needs: where the problem was actually seen. An inferred fork
baseline, for example from old tickets loaded into HPB, is flagged as
such.

\paragraph{Steps.}
What happens inside a branch is a short, time-ordered list of
\emph{steps}: done, to do, test passed, test failed, each optionally with
evidence such as a report, a measurement or a build identifier. Steps are
where laboratory results and verifications live, and they are what the
release gate checks.

\subsection{How a branch lives and ends}
\label{sec:model-closing}

A branch is in one of four states (\cref{def:state},
\cref{fig:lifecycle}):
\begin{itemize}[leftmargin=1.5em,itemsep=1pt,topsep=2pt]
  \item \textsf{open}: accepted, not yet committed to a version;
  \item \textsf{targeted}: committed to a planned baseline, its
        \emph{target} $g(b)$. The target can be changed (re-target) or
        withdrawn (un-target), and the change remains visible;
  \item \textsf{merged}: its result is part of a released baseline,
        $m(b)$;
  \item \textsf{discarded}: closed without a baseline, with a written
        rationale.
\end{itemize}
Merged and discarded are final. Reopening a closed branch exists, but only
as an exceptional, audited operation.

% F7 -- Life cycle of a branch: the state machine of Definition "Branch state".
% Requires hpb.sty v1.0 (palette and glyph styles). Plain TikZ, units in cm.
\begin{figure}[t]
\centering
\resizebox{\linewidth}{!}{%
\begin{tikzpicture}[
    x=1cm, y=1cm,
    font=\sffamily\fontsize{7}{8.4}\selectfont,
    lc/state/.style   = {rectangle, rounded corners=5pt, draw=hpbInk,
                         line width=1pt, minimum width=3.3cm,
                         minimum height=1.45cm, inner sep=0pt},
    lc/final/.style   = {double=white, double distance=1.3pt, line width=0.8pt},
    lc/name/.style    = {font=\sffamily\bfseries\fontsize{8}{9.6}\selectfont,
                         text=hpbInk, anchor=west, inner sep=0pt},
    lc/inv/.style     = {font=\sffamily\fontsize{6.5}{7.8}\selectfont,
                         text=hpbMuted, anchor=west, inner sep=0pt},
    lc/arrow/.style   = {draw=hpbInk, line width=1pt, rounded corners=6pt,
                         -{Stealth[length=5pt,width=4.5pt]}},
    lc/audit/.style   = {lc/arrow, dash pattern=on 4pt off 2.5pt},
    lc/tr/.style      = {font=\sffamily\fontsize{7}{8.4}\selectfont,
                         text=hpbInk, fill=white, inner sep=1.5pt},
    lc/trb/.style     = {lc/tr, font=\sffamily\bfseries\fontsize{7}{8.4}\selectfont},
  ]
  % ---- the open set W(t): states in which a branch counts as work in progress
  % (dotted frame: deliberately not the style of an HPB gate band)
  \draw[draw=hpbDiscInk, line width=0.8pt, rounded corners=4pt,
        dash pattern=on 0.1pt off 2pt, line cap=round]
    (1.45,-1.02) rectangle (10.15,1.02);
  \node[lc/inv, anchor=south west, font=\sffamily\itshape\fontsize{6.5}{7.8}\selectfont]
    at (1.5,1.06) {open set $W(t)$};

  % ---- states: neutral fills (not the baseline color of the notation);
  %      the glyph in each box is the HPB terminal that represents the state
  \node[lc/state, fill=white]               (open) at (3.35,0)  {};
  \node[lc/state, fill=white]               (targ) at (8.25,0)  {};
  \node[lc/state, lc/final, fill=white]     (merg) at (13.65,0) {};
  \node[lc/state, lc/final, fill=white]     (disc) at (8.25,-3) {};

  % glyphs (shown for an improvement)
  \node[hpb/head ring, draw=hpbImp] at (2.2,0) {};
  \node[hpb/event, fill=white]      at (2.2,0) {};
  \draw[hpb/path, hpb/path target] (7.3,0) -- (7.83,0);
  \node[hpb/head ring, draw=hpbImp] at (7.1,0) {};
  \node[hpb/event, fill=white]      at (7.1,0) {};
  \node[hpb/event, fill=hpbImp]     at (12.5,0) {};
  \draw[draw=hpbInk, line width=1.3pt, line cap=round, line join=round]
    (12.5,0) ++(-2.6pt,0.1pt) -- ++(1.8pt,-2pt) -- ++(3.5pt,4.2pt);
  \node[hpb/event, fill=hpbDisc, draw=hpbMuted] at (7.1,-3) {};
  \draw[draw=hpbInk, line width=1.3pt, line cap=round]
    (7.1,-3) ++(-2.2pt,-2.2pt) -- ++(4.4pt,4.4pt) ++(-4.4pt,0) -- ++(4.4pt,-4.4pt);

  % names and invariants
  \node[lc/name] at (3.0,0.22)   {Open};
  \node[lc/inv]  at (3.0,-0.2)   {$g(b)$ unset};
  \node[lc/name] at (7.95,0.22)  {Targeted};
  \node[lc/inv]  at (7.95,-0.2)  {$g(b)$ planned;};
  \node[lc/inv]  at (7.95,-0.47) {verified at $t_r(b)$};
  \node[lc/name] at (13.4,0.22)  {Merged};
  \node[lc/inv]  at (13.4,-0.2)  {$m(b)=g(b)$};
  \node[lc/name] at (7.95,0.22-3) {Discarded};
  \node[lc/inv]  at (7.95,-0.2-3) {rationale $\rho(b)$};

  % ---- creation: a branch is opened by forking from a released baseline,
  %      drawn as the released-baseline glyph of the notation
  \node[hpb/baseline] (init) at (-0.35,0) {};
  \draw[lc/arrow] (init.east) -- (open.west)
    node[lc/trb, midway, above=2pt] {fork}
    node[lc/inv, midway, below=3pt, anchor=north] {from $f(b)$};

  % ---- ordinary transitions
  \draw[lc/arrow] ([yshift=3.5mm]open.east) -- ([yshift=3.5mm]targ.west)
    node[lc/trb, midway, above=1pt] {target};
  \draw[lc/arrow] ([yshift=-3.5mm]targ.west) -- ([yshift=-3.5mm]open.east)
    node[lc/trb, midway, below=1pt] {un-target};
  \draw[lc/arrow] (7.75,0.725) -- (7.75,1.3) -- (8.75,1.3) -- (8.75,0.725);
  \node[lc/trb, anchor=south] at (8.25,1.32) {re-target};
  \draw[lc/arrow] (targ.east) -- (merg.west)
    node[lc/trb, midway, above=2pt] {merge}
    node[lc/inv, midway, below=3pt, anchor=north, align=center]
      {on publication\\of $g(b)$, if verified};
  \draw[lc/arrow] ([xshift=6mm]open.south) |- (disc.west)
    node[lc/trb, pos=0.75, above=1pt] {discard};
  \draw[lc/arrow] (targ.south) -- (disc.north)
    node[lc/trb, midway, right=2pt] {discard};

  % ---- exceptional, audited transitions (dashed)
  \draw[lc/audit] (merg.north) -- (13.65,1.95) -- (3.95,1.95) -- (3.95,0.725);
  \node[lc/trb, anchor=south] at (8.8,1.97) {reopen\quad
    \textmd{\textcolor{hpbMuted}{(exceptional, audited)}}};
  \draw[lc/audit] (disc.south) -- (8.25,-4.15) -- (2.75,-4.15) -- (2.75,-0.725);
  \node[lc/trb, anchor=north] at (5.5,-4.17) {reopen\quad
    \textmd{\textcolor{hpbMuted}{(exceptional, audited)}}};
\end{tikzpicture}}
\caption{Life cycle of a branch (\cref{def:state}). A branch is
\textsf{open} until it is committed to a planned version (\textsf{targeted});
it merges only from \textsf{targeted}, when that version is published and
after it has been verified; it can be discarded from either state, always
with a reason. \textsf{Merged} and \textsf{discarded} are final; dashed
transitions are exceptional and audited. Each box shows the glyph of the
state; the dotted frame marks the states counted as work in progress,
$W(t)$.}
\label{fig:lifecycle}
\end{figure}


\paragraph{Merge.}
A branch is merged only from \textsf{targeted}, and only when its target
baseline is published (\cref{def:merge}). The relation ``branch $b$ is
closed by baseline $v$'' is what produces the change log of $v$. This
looks strict, and it is meant to be: every merge fulfills a commitment that
was written down beforehand, and released versions are never altered by
ordinary work. For an urgent fix, targeting, verification and publication
can follow each other within hours.

A targeted branch is \emph{verified} when it has a passed test step whose
evidence identifies a candidate of its target: a release-candidate build,
a prototype or a pilot lot. Verification is checked at the gate.

\paragraph{Discard.}
A discard closes a branch without a baseline and always carries a
rationale from one of two classes. \emph{Abandoned}: not viable,
duplicate, cannot reproduce, superseded, out of scope, will not fix.
\emph{Resolved without a product baseline}: already resolved in an
existing version, contained without design change (for example by
reworking stock), resolved by the supplier. Keeping the two classes apart
matters for quality reporting: a problem solved outside the product is not
the same as a problem nobody solved. A defect whose solution changes
\emph{how} the product is built rather than its design is not discarded;
it is merged into a process baseline.

\subsection{Line forks, gates and the life of a reference}
\label{sec:model-forks}

\paragraph{Line fork.}
A \emph{line fork} is a hardware baseline that creates a new reference
$P'$ (\cref{def:linefork}). The hardware baseline is the first baseline of
the new line, and the new line remembers its parent and the parent
baseline it was derived from. The branches that produce the new reference,
typically lab evolutions and improvements, fork from the parent line and
merge into that hardware baseline; later work forks from the new line.
Two consequences are worth stating because they are exactly the questions
that come up in practice. First, later versions of the parent do
\emph{not} contain the work that created the new reference. Second, a fix
made on the parent after the fork is \emph{not} in the child unless it is
carried over explicitly (\cref{fig:hardware-fork}).

\paragraph{Release gate.}
Every planned baseline is a \emph{gate} (\cref{def:gate}). Its
\emph{targeted set} $\mathcal{B}(g)$ is the set of branches currently
committed to it. Publishing the baseline is a single operation that
releases it and merges every branch of its targeted set. Before that,
every member must be verified, and anything it depends on must already be
released or be released in the same operation; whatever is not ready is
re-targeted or discarded first (\cref{rule:gate}). This is the
Stage-Gate idea~\citep{cooper1990stagegate} applied to versions instead of
projects, and it removes the most common source of misleading release
notes: work that was ``in'' the version on the plan but quietly did not
make it (\cref{fig:release-gate}).

\paragraph{Dependencies, propagations and launch bundles.}
Three optional relations cover the cases where versions and branches of
different lines are tied together (\cref{def:relations}). A
\emph{dependency} says that a branch can only be delivered with or after
another branch or baseline, for example a firmware improvement that needs
hardware \texttt{Rev~B}, or work blocked by a supplier sample. A
\emph{propagation} says that a branch carries the change of another branch
to a different stream or line: a back-port or forward-port. A change that
must reach several versions is therefore one branch per target, linked by
propagations, never one branch with several endings (\cref{rule:closure});
that is how a missing back-port becomes visible. A \emph{launch bundle}
is a set of planned baselines of different lines that are published
together or not at all, for example hardware \texttt{Rev~B} with firmware
\texttt{v2.0} and app \texttt{5.0}.

\paragraph{Life cycle of a reference.}
In a portfolio with high rotation, references are retired as often as
they are created, and the rules for what may still change on an old
reference are usually tribal knowledge. HPB makes them explicit. Every
line has up to four dated milestones: end of sale, end of production,
start of security-only maintenance and end of service
(\cref{def:lifecycle}). They divide its life into phases, and
\cref{tab:phases} in \cref{app:formal} says which branches and baselines
each phase still accepts (\cref{rule:lineage}). Supplier notices are a frequent trigger in
the later phases: a product change notice or an end-of-life notice for a
component usually opens a lab evolution (qualify an alternative) or a task
(last-time buy), targeted at a hardware baseline that is either
interchangeable or, if form, fit or function change, a line fork.


\paragraph{Deployment overlay.}
The questions a product manager is asked most often are about customers:
\emph{who is exposed to this defect, who has not upgraded, which promise
has slipped?} An optional \emph{deployment overlay} (\cref{def:deployment})
records the rollout status and installed base of each released baseline,
the configuration each customer runs, and the commitments made on each
branch. ``Customers exposed to X-118'' is then simply the set of customers
whose configuration contains a version in which X-118 occurs.

\subsection{What you get for free}
\label{sec:model-free}

The point of the rules is that the documents a product manager usually
compiles by hand become \emph{queries on the graph}. Four of them are
stated precisely, with proof sketches, in \cref{app:formal}:
\begin{itemize}[leftmargin=1.5em,itemsep=1pt,topsep=2pt]
  \item \textbf{Change log.} The content of a version is the set of
        branches it closed; grouped by type, it is the release note
        (\cref{prop:changelog}).
  \item \textbf{Escaped defects.} Every issue belongs to exactly one
        version: the one in which it was observed (\cref{prop:escaped}).
  \item \textbf{Work in progress.} At any moment, the open and targeted
        branches are a well-defined set, split by type, to which Little's
        law~\citep{little1961} applies (\cref{prop:wip}).
  \item \textbf{Lineage.} Every reference has a single ancestry and every
        version a single chain of predecessors, so ``is this fix in that
        version?'' always has a yes-or-no answer (\cref{prop:lineage}).
\end{itemize}
The last point has a practical corollary that I consider one of the most
useful things HPB shows. A fix merged on a maintenance stream, or on a
parent line after a line fork, is not contained in the main stream or the
child line. The diagram makes the question ``must this fix be carried
over?'' visible, and the rules force the answer to be recorded: either as
a propagated branch or as an explicit decision (\cref{fig:backport}).
